\section{实现局限与改进方向}
距离交会的法方程迭代没有按测距精度加权，也没有剔除 NLOS 离群值；改进时应采用基于测距方差的加权最小二乘、鲁棒损失、几何精度因子与残差诊断，并通过独立标定确认基站坐标和设备时间同步。RSS 定位可增加 AP 方差建模、缺失机制处理、设备 RSS 偏差校准以及按轨迹段划分的留出验证；按时间随机拆分会产生空间/时间泄漏风险，应优先采用独立采集段验证。PDR 可标定步长与陀螺零偏，改用自适应步长和磁干扰检测，并在有可信锚点时融合几何或指纹观测来约束漂移。

本次工作范围还有明确边界：除指纹算法外，没有重新执行所有原有 C++ 变体；无人机、UWB、学生无遮挡/遮挡实验未形成新的独立复算结果。原始输出文件的来源版本也未由本轮 Git 历史逐项确认。因此报告将教室/PDR结果称为既有程序输出的复核，不作“本轮完整重跑”的表述。

\section{结论}
本项目形成了从测距、RSS 匹配到惯性航迹推算的互补实验链。指纹算法在给定测试集上实现了 100\% 输出率，水平误差 RMSE 为 1.422 m；教室超声既有解算经参考轨迹复核得到 0.239 m RMSE，测距拟合残差中位数约 9.9 mm；PDR 既有轨迹按步伐时刻与经纬度参考路径对照得到 8.381 m RMSE。测距拟合残差与坐标误差含义不同，而跨实验数据条件也限制了算法横向排名。最值得进一步完成的工作是隔离重跑全部 C++ 程序、统一坐标系/时间基准，并基于独立重复采集开展参数标定和跨场景验证。

\section*{参考资料}
\addcontentsline{toc}{section}{参考资料}
\begin{enumerate}[label={[\arabic*]},leftmargin=2.2em]
\item 室内定位课程课件：\textit{室内定位技术——几何定位}、\textit{室内定位技术——航迹推算}、\textit{室内定位技术——数据库匹配}、\textit{室内定位传感器}，课程提供文件，2026。
\item IndoorPos 项目源码：\url{src/DistanceIntersection/}、\url{src/Fingerprint/}、\url{src/PDR/}、\url{src/Sensors/}，课程实验材料。
\item IndoorPos 实验数据说明：距离交会数据目录下 \texttt{数据格式.md/readme.md}、指纹定位数据及 PDR 数据文件，课程实验材料。
\item Android Sensors 数据采集说明：配套使用文档及 \texttt{DataConvert} MATLAB 转换程序，项目配套材料。
\end{enumerate}
